\par\noindent\begin{minipage}{\linewidth}
% !TeX root = ../教材题库_全部条目_转写工作稿.tex
% Source: XJ-B1-C03-PR-P076-Y0567
% 第3章 / 3.1 函数 / 教材 p.76 / PDF 第83页
% 已对照原页转写题干及公式；待教师复核；未制作答案
% 原题及相邻原解截图（核对用）：../题目原图/XJ-B1-C03-PR-P076-Y0567-01.webp
% 本文件只在转写完成后作为题干使用；OCR 条目尚不是可靠题干。
\begin{bankitem}{p.76 练习 第3题}{第3章 / 3.1 函数 / 教材 p.76 / PDF 第83页}{XJ-B1-C03-PR-P076-Y0567}{已对照原页转写题干及公式；待教师复核；未制作答案}
一个质点沿直线运动。质点由静止匀加速 $T\,\mathrm s$ 后速度达到 $8\,\mathrm{m/s}$，然后质点以恒定速度 $8\,\mathrm{m/s}$ 运动了 $5T\,\mathrm s$，之后质点在 $40\,\mathrm s$ 内匀减速到完全停下。
\subq{1}画出质点运动的速度-时间图象；
\subq{2}已知质点运动的总位移是 $600\,\mathrm m$，求 $T$ 的值；
\subq{3}画出质点运动的加速度-时间图象。
\end{bankitem}
\end{minipage}\par\medskip
